% META: {"id":"1SPE-VARALEA-ME-003","chapitre":"1SPE-VARIABLES-ALEATOIRES","type_objet":"methode","capacites_codes":[],"extension_codes":["X1"],"programme_alignment":"OPTIONAL_EXTENSION","extension_label":"Approfondissement — Vers la Terminale","methodes":["M3"],"sources_inspiration":[],"status":"approved","origin":"generated","status_history":["generated","audited","accepted","approved"],"release_acceptance":"RELEASE_OWNER_BATCH_ACCEPTANCE","acceptance_closure_digest":"sha256:9b3ccf9a81c5520fbb7e03b7b2d3e7bf2057a02834b6d908bf3be5f49f3b553f","acceptance_receipt_digest":"sha256:4fad86f0282e0fa4e0419cca1ad6379ae7af5a280c04a4a90880f90a1c044242"}
\begin{center}\textbf{Approfondissement — Vers la Terminale}\end{center}
\begin{fichemethode}{M3}{Reconnaître et utiliser une loi binomiale}

\quandUtiliser{%
  Utiliser cette méthode lorsque l'on répète une même expérience à deux issues (succès/échec) de façon indépendante et que l'on compte le nombre de succès.
}

\pasApas{%
  \begin{enumerate}
    \item Je vérifie les trois conditions du schéma de Bernoulli :
      \begin{itemize}
        \item Deux issues seulement (succès de probabilité $p$, échec de probabilité $1-p$).
        \item Les épreuves sont \textbf{identiques} (même $p$ à chaque fois).
        \item Les épreuves sont \textbf{indépendantes}.
      \end{itemize}
    \item J'identifie $n$ (nombre de répétitions) et $p$ (probabilité du succès).
    \item Je pose $X$ le nombre de succès : $X \sim \mathcal{B}(n, p)$.
    \item Pour calculer $P(X = k)$, j'applique : $P(X = k) = \binom{n}{k} p^k (1-p)^{n-k}$.
    \item Je calcule le coefficient binomial $\binom{n}{k}$ (triangle de Pascal ou formule).
  \end{enumerate}
}

\exempleRedige{%
  On lance $5$ fois une pièce équilibrée. Quelle est la probabilité d'obtenir exactement $3$ Pile ?

  \medskip
  \commentaireMarge{Vérifier les conditions.}%
  Succès = Pile, $p = \frac{1}{2}$, $n = 5$ lancers indépendants.

  \commentaireMarge{$X \sim \mathcal{B}(5, \frac{1}{2})$.}%
  $X$ = nombre de Pile.

  \commentaireMarge{Calcul de $P(X = 3)$.}%
  $P(X = 3) = \binom{5}{3} \left(\frac{1}{2}\right)^3 \left(\frac{1}{2}\right)^2 = 10 \times \frac{1}{8} \times \frac{1}{4} = \frac{10}{32} = \frac{5}{16}$.

  \commentaireMarge{Vérification rapide.}%
  $\binom{5}{3} = \frac{5!}{3!2!} = \frac{120}{6 \times 2} = 10$ \checkmark. Somme des exposants : $3 + 2 = 5$ \checkmark.
}

\pieges{%
  \begin{itemize}
    \item \textbf{Piège 1.} Appliquer la loi binomiale sans vérifier l'indépendance (tirage sans remise, par exemple).
    \item \textbf{Piège 2.} Inverser les exposants : $p^k$ (succès) et $(1-p)^{n-k}$ (échecs), pas l'inverse.
    \item \textbf{Piège 3.} Oublier le coefficient binomial $\binom{n}{k}$ dans la formule.
  \end{itemize}
}

\verifier{%
  La somme des exposants de $p$ et de $(1-p)$ doit être $n$. Le coefficient binomial $\binom{n}{k}$ est un entier positif.
}

\sEntrainer{\refExos{M3}}

\end{fichemethode}
